% ================================================
% 文件: 04_线材.tex
% 章节: 第5章 线材
% 所属部分: 电子元器件与材料
% 版本: v1.0
% ================================================
% 本章目录结构（树状）:
% \chapter{线材}
% ├── \section{线材基础}
% ├── \section{导线类型与特性}
% ├── \section{电缆与传输线}
% ├── \section{连接器与接头}
% └── \section{线材选型与应用}
% ================================================

\chapter{线材}
\label{chap:wiring_materials}

\section{线材基础}
\label{sec:wiring_fundamentals}

线材是电子设备中用于传输电能和信号的导体材料。线材的主要参数包括：

\begin{itemize}
    \item \textbf{导体材料}：铜、铝、银等，铜是最常用的导体材料
    \item \textbf{截面积}：常用AWG（美国线规）或mm²表示
    \item \textbf{绝缘材料}：PVC、PE、PTFE等
    \item \textbf{额定电压}：导线安全工作的最高电压
\end{itemize}

线材虽是设备中最不引人注目的部分，却承担着两项互不相同的任务：
一是把电能从电源输送到各级负载，此时关注的是载流能力、压降与
发热；二是把信息从一处传送到另一处，此时关注的是特性阻抗、
衰减、串扰与屏蔽效能。前者属于电力配线问题，后者属于传输线
问题。频率越高、线路越长，第二类考虑就越占主导地位。通信设备
故障统计表明，相当高比例的间歇性故障来源于导线断股、绝缘老化
和接头氧化，因此对线材的正确认识与规范施工，其重要性并不低于
对有源器件的选择。

\textbf{导体材料}：衡量导体优劣的首要指标是电阻率 $\rho$。
在常用金属中，银的电阻率最低，约
$1.59\times10^{-8}\ \Omega\cdot\mathrm{m}$；铜次之，约
$1.72\times10^{-8}\ \Omega\cdot\mathrm{m}$；铝约为
$2.83\times10^{-8}\ \Omega\cdot\mathrm{m}$，即约为铜的 1.64 倍。
银的导电性虽最好，但价格昂贵，通常只用于表面镀层，以改善高频
导体的表面导电性与接触点的抗氧化性能。铝的电阻率较高，同样载流
能力下截面必须加大约 $60\%$，但密度仅为铜的三分之一，故在长距离
架空电力线路中仍具优势；其缺点是易生成致密氧化膜、蠕变大、
不易可靠焊接，因此设备内部配线极少采用。综合导电性、机械强度、
可焊性与成本，铜是最常用的导体材料。设备用铜线多为软态退火铜
（TU、T2 等牌号），其伸长率大、柔韧性好；导体表面常镀锡以
便于焊接并防止硫化发黑，射频导体则镀银以降低表面损耗。

\textbf{截面积与线规}：导体的截面积直接决定其直流电阻与载流
能力。我国与欧洲习惯以平方毫米（mm²）标注公称截面，如
$0.5\ \mathrm{mm}^2$、$1.0\ \mathrm{mm}^2$、$2.5\ \mathrm{mm}^2$；
北美则采用 AWG（美国线规）编号，编号越大导线越细。AWG 的
线径按几何级数递减，其直径与线号的关系可表示为
\begin{equation}
d = 0.127 \times 92^{\,(36-n)/39} \ \ (\mathrm{mm})
\end{equation}
式中 $n$ 为 AWG 线号。由该式可推出两条便于记忆的规律：线号
每减小 3，截面积约增大一倍；线号每减小 6，直径约增大一倍、
截面积约增大四倍。例如 AWG 10 的截面约
$5.26\ \mathrm{mm}^2$，AWG 16 约 $1.31\ \mathrm{mm}^2$，
AWG 22 约 $0.326\ \mathrm{mm}^2$。需注意多股线的标称截面是
各股导体截面之和，其外径因绞合间隙而略大于同截面的单芯线。

\textbf{绝缘材料}：绝缘层的作用是防止导体之间及导体对地短路，
同时提供机械保护。常用材料包括聚氯乙烯（PVC）、聚乙烯（PE）、
交联聚乙烯（XLPE）、聚四氟乙烯（PTFE，俗称铁氟龙）、
硅橡胶与聚酰亚胺。PVC 成本低、加工方便、耐油耐酸，长期工作
温度一般为 $70\ ^{\circ}\mathrm{C}$（耐热型可达
$105\ ^{\circ}\mathrm{C}$），但介电损耗较大且燃烧时释放
有害气体，仅适用于低频与电力配线。PE 的相对介电常数约 2.3、
损耗角很小，是同轴电缆与网线的主要绝缘材料；发泡处理后
介电常数可降至 1.5 左右，能显著降低高频衰减并提高速度因数。
PTFE 耐温可达 $200\ ^{\circ}\mathrm{C}$ 以上，介电常数约 2.1
且随频率与温度变化极小，耐化学腐蚀，是射频同轴电缆与航空
航天线材的首选，缺点是价格高、质地较硬。硅橡胶柔软耐寒耐热，
用于需反复弯折或高温环境的引线。

\textbf{额定电压与温度等级}：额定电压是导线安全工作的最高
电压，取决于绝缘层材料与厚度，设备内部配线常见等级为
$300\ \mathrm{V}$、$600\ \mathrm{V}$ 与
$1000\ \mathrm{V}$。选用时应使工作电压不超过额定值，且对
脉冲与浪涌留有余量。与额定电压同等重要的是温度等级，它规定
了绝缘层长期允许的工作温度；导线的实际温升由自身铜损、
环境温度与散热条件共同决定，成束敷设或穿管时散热恶化，
必须相应降低载流量。绝缘寿命对温度极为敏感，经验上温度每
升高约 $10\ ^{\circ}\mathrm{C}$，绝缘老化速率约加倍，
这就是电力配线必须留有充分裕量的根本原因。

\section{导线类型与特性}
\label{sec:wire_types}

常见的导线类型包括：

\begin{itemize}
    \item \textbf{单芯线}：单根导体，用于固定安装
    \item \textbf{多股线}：多根细导体绞合，柔韧性好
    \item \textbf{漆包线}：表面涂漆绝缘，用于绕制线圈
    \item \textbf{电磁线}：用于电机、变压器绕组
\end{itemize}

导线的直流电阻：$R = \rho\frac{l}{S}$，其中$\rho$为电阻率，$l$为长度，$S$为截面积。

单芯线（实心线）由一根导体构成，外径小、结构稳定、便于压接
与穿管，适用于一次装配后不再移动的固定安装场合，如印制板
跳线、配线架内部走线与建筑物内的固定敷设。其缺点是反复弯折
易在弯曲处产生疲劳裂纹而折断。多股线由多根细导体绞合而成，
在相同截面下柔韧性显著提高，弯折寿命长，适用于活动部位、
设备之间的连接线与手持设备引线；股数越多、单股越细则越柔软，
如常见的 $7$ 股、$19$ 股、$41$ 股结构。多股线在焊接时应先
上锡固化，但不宜把上锡段压入螺钉端子，因为锡在长期压力下
会发生蠕变而使接触压力下降。

漆包线的导体表面涂覆聚酯、聚氨酯或聚酰亚胺漆膜作为绝缘，
漆膜极薄，故绕制线圈时槽满率高、体积小，是电感器、
中频变压器、继电器线圈的标准用线。按耐温等级分为
$130\ ^{\circ}\mathrm{C}$（B 级）、$155\ ^{\circ}\mathrm{C}$
（F 级）、$180\ ^{\circ}\mathrm{C}$（H 级）等。焊接时须先
刮除或用高温焊锡烧穿漆膜，自粘性与直焊性漆包线可省去这一
工序。电磁线是漆包线、丝包线、纸包线与扁绕线的统称，专用于
电机、变压器与电磁铁绕组，除电气性能外还要求良好的机械
强度与耐浸渍性。用于高频场合的利兹线（多股细漆包线绞合并
分层换位）能有效削弱趋肤效应与邻近效应，从而提高线圈的
$Q$ 值。

导线的直流电阻按 $R = \rho l/S$ 计算，并随温度变化。以
$20\ ^{\circ}\mathrm{C}$ 为基准，铜导线的电阻温度关系为
\begin{equation}
R_t = R_{20}\left[\,1 + \alpha\,(t - 20)\,\right], \qquad
\alpha_{\mathrm{Cu}} \approx 0.00393\ /^{\circ}\mathrm{C}
\end{equation}
即温度由 $20\ ^{\circ}\mathrm{C}$ 升至
$70\ ^{\circ}\mathrm{C}$ 时电阻约增大 $20\%$。这一效应在
计算满载压降与保护整定值时不可忽略。导线上的电压降为
$\Delta U = I R$，对于往返两根导线的供电回路应按双倍长度
计算，即 $\Delta U = 2 I \rho l / S$；相应的功率损耗为
$P = I^2 R$，全部转化为热量。工程上通常要求直流供电回路的
线路压降不超过额定电压的 $3\%\sim5\%$，对低压大电流回路
（如 $13.8\ \mathrm{V}$ 电台电源线）这一限制往往比载流量
更为苛刻，是决定线径的实际依据。

\begin{table}[htbp]
\centering
\caption{常用线规、截面积、铜导线电阻与参考载流量}
\begin{tabular}{|p{1.8cm}|p{2.2cm}|p{2.0cm}|p{2.4cm}|p{2.4cm}|}
\hline
\textbf{AWG} & \textbf{截面积/mm$^2$} & \textbf{直径/mm} & \textbf{电阻/($\Omega$/100m)} & \textbf{参考载流量/A} \\
\hline
10 & 5.26 & 2.588 & 0.33 & 25$\sim$30 \\
\hline
12 & 3.31 & 2.053 & 0.52 & 18$\sim$20 \\
\hline
14 & 2.08 & 1.628 & 0.83 & 12$\sim$15 \\
\hline
16 & 1.31 & 1.291 & 1.31 & 8$\sim$10 \\
\hline
18 & 0.823 & 1.024 & 2.09 & 5$\sim$7 \\
\hline
20 & 0.518 & 0.812 & 3.32 & 3$\sim$5 \\
\hline
22 & 0.326 & 0.644 & 5.28 & 2$\sim$3 \\
\hline
24 & 0.205 & 0.511 & 8.39 & 1.5$\sim$2 \\
\hline
26 & 0.129 & 0.405 & 13.3 & 1.0$\sim$1.3 \\
\hline
\end{tabular}
\end{table}

表中载流量仅为设备内部单根明敷、环境温度不高于
$30\ ^{\circ}\mathrm{C}$ 时的参考值。若多根成束敷设、
穿入线管或环境温度较高，必须按相应系数降额；反之，短距离
的机内跳线可适当放宽。实际工程应以所选电缆制造厂给出的
数据为准。

\section{电缆与传输线}
\label{sec:cables_transmission_lines}

电缆是由多根导线组成的传输线路，常见类型：

\begin{itemize}
    \item \textbf{同轴电缆}：内导体、绝缘层、外导体、护套，用于高频信号传输
    \item \textbf{双绞线}：两根导线绞合，用于差分信号传输
    \item \textbf{屏蔽电缆}：带有屏蔽层，抗干扰能力强
    \item \textbf{电力电缆}：用于传输大功率电能
\end{itemize}

同轴电缆的特性阻抗通常为50Ω或75Ω。

当信号波长与线路长度可以相比时，导线不再是简单的连接件，
而必须视为具有分布参数的传输线。传输线由分布电感 $L_0$、
分布电容 $C_0$、分布电阻与漏电导构成，在损耗可忽略的条件下
其特性阻抗为
\begin{equation}
Z_0 = \sqrt{\frac{L_0}{C_0}}
\end{equation}
特性阻抗只取决于导体的几何尺寸与介质，而与线的长度无关。
对同轴结构，若外导体内径为 $D$、内导体外径为 $d$、绝缘介质
相对介电常数为 $\varepsilon_r$，则
\begin{equation}
Z_0 = \frac{138}{\sqrt{\varepsilon_r}}\,\lg\frac{D}{d}\ \ (\Omega)
\end{equation}
信号在电缆中的传播速度低于真空光速，其比值称为速度因数
$v/c = 1/\sqrt{\varepsilon_r}$：实心聚乙烯绝缘约为 $0.66$，
发泡聚乙烯约为 $0.80$，聚四氟乙烯约为 $0.70$。这一数值在
制作四分之一波长匹配线与延迟线时必须准确使用。

同轴电缆由内导体、绝缘层、外导体（编织网或铝箔）与外护套
组成，其电磁场被完全约束在内外导体之间，故辐射与拾取干扰
都很小，是用于高频信号传输的主要形式。射频系统的标准
特性阻抗为 $50\ \Omega$，这是在给定外径下兼顾功率容量与
衰减而折中的结果；有线电视与视频系统采用 $75\ \Omega$，
因该值使同轴结构的衰减最小。当传输线两端阻抗与 $Z_0$
不相等时会产生反射，反射系数与电压驻波比分别为
\[
\Gamma = \frac{Z_L - Z_0}{Z_L + Z_0}, \qquad
\mathrm{VSWR} = \frac{1+|\Gamma|}{1-|\Gamma|}
\]
驻波比过高不仅使发射功率无法有效送出，还会在馈线上形成
电压电流的极大点，可能引起绝缘击穿或功放管损坏，因此
天馈系统调试的核心指标之一就是把驻波比压低到 $1.5$ 以下。
电缆的衰减以 dB/100m 表示，随频率升高而增大（低频段主要
受导体损耗支配，近似与 $\sqrt{f}$ 成正比；高频段介质损耗
逐渐显著，与 $f$ 近似成正比），且粗电缆的衰减小于细电缆。

双绞线把两根导线以一定节距绞合，使外界磁场在相邻绞距内
产生的感应电动势相互抵消，从而大幅降低共模干扰的耦合；
配合差分收发电路，可在不加屏蔽的条件下获得良好的抗扰性，
这是它被广泛用于差分信号传输的原因。其特性阻抗典型值为
$100\ \Omega$（局域网五类线）或 $120\ \Omega$（现场总线），
按有无屏蔽层分为 UTP 与 STP。屏蔽电缆在绝缘层外增加铝箔、
铜编织网或双层复合屏蔽，屏蔽层须可靠接地才能发挥作用；
低频时屏蔽层宜单端接地以避免地环流，射频时则应两端
（乃至沿途多点）接地以维持屏蔽连续性。电力电缆导体截面大、
绝缘厚、机械保护强，用于传输大功率电能，选型时以载流量、
短路耐受与压降为主要依据。

\begin{table}[htbp]
\centering
\caption{几种常用同轴电缆与传输线的特性}
\begin{tabular}{|p{2.4cm}|p{2.0cm}|p{2.2cm}|p{4.2cm}|}
\hline
\textbf{型号/类别} & \textbf{阻抗} & \textbf{速度因数} & \textbf{特点与典型用途} \\
\hline
RG-174 & 50 $\Omega$ & 约 0.66 & 极细软，衰减大；仪表跳线、机内短连接 \\
\hline
RG-58 & 50 $\Omega$ & 约 0.66 & 常用细馈线，短距离收发馈电与测试线 \\
\hline
RG-213 & 50 $\Omega$ & 约 0.66 & 粗馈线，衰减小、功率容量大；固定台馈线 \\
\hline
RG-59 & 75 $\Omega$ & 约 0.66 & 视频与射频分配，有线电视入户线 \\
\hline
五类双绞线 & 100 $\Omega$ & 约 0.65 & 差分传输，以太网与数据链路 \\
\hline
平行馈线 & 300$\sim$600 $\Omega$ & 约 0.9 & 损耗极低，需平衡馈电与远离金属体 \\
\hline
\end{tabular}
\end{table}

\section{连接器与接头}
\label{sec:connectors}

连接器用于连接导线或设备，常见类型：

\begin{itemize}
    \item \textbf{RF连接器}：SMA、BNC、N型等，用于射频信号
    \item \textbf{音频连接器}：XLR、RCA、3.5mm耳机插头
    \item \textbf{电源连接器}：插头插座、端子排
    \item \textbf{电路板连接器}：排针排母、板对板连接器
\end{itemize}

连接器的关键参数包括接触电阻、绝缘电阻、插拔次数等。

连接器是线路中最薄弱的环节：它既要保证可靠的电接触，又要
允许反复分离，还要在振动、温变与腐蚀条件下长期维持性能。
对射频连接器而言，另有一项苛刻要求——必须保持传输线的
阻抗连续性。SMA 采用螺纹连接，尺寸小、性能好，可用至
$18\ \mathrm{GHz}$ 甚至更高，适于仪表与机内互连，但插拔
寿命有限且拧紧力矩需符合规定。BNC 采用卡口式连接，插拔
迅速、便于频繁操作，常用频率在 $4\ \mathrm{GHz}$ 以下，
广泛用于示波器探头、视频与中频信号。N 型连接器体积较大、
功率容量高、密封性好，是基站与固定台馈线的常用接口。
此外还有 UHF 型（早期通用射频接口，阻抗不连续、仅适用于
较低频率）、TNC（带螺纹的 BNC 改进型，抗振性更好）与
7/16 DIN（大功率低互调场合）。选用时须注意阻抗规格：
同一外形的连接器可能有 $50\ \Omega$ 与 $75\ \Omega$ 两种
版本，混用会在接头处引入反射。

音频连接器中，XLR 为三芯卡侬接口，用于平衡传输，其两根
信号线加屏蔽的结构可在长距离传输中有效抑制共模噪声，是
专业音频与广播设备的标准；RCA（莲花插头）为非平衡接口，
用于消费类音视频设备；$3.5\ \mathrm{mm}$ 耳机插头结构
紧凑，按环节数量分为二芯（单声道）、三芯（立体声或
单声道带话筒）与四芯（立体声带话筒），使用时须注意各
环节定义在不同厂家间可能不同。通信设备的手咪、耳机接口
多采用多芯圆形航空插头，具有锁紧螺套与防误插定位键。

电源连接器与端子排的首要指标是额定电流与接触压力。压接
端子必须使用与线径匹配的专用压线钳，压接后既不可出现
毛刺切断线股，也不可松动；螺钉端子应按规定扭矩紧固，
并定期复检，因为热循环会使接触压力逐渐松弛。电路板
连接器包括排针排母、FFC/FPC 扁平电缆座、板对板连接器与
卡边连接器，其间距已由 $2.54\ \mathrm{mm}$ 逐步减小到
$1.27\ \mathrm{mm}$、$0.5\ \mathrm{mm}$ 以下，装配时须
注意防止错位与虚接。

连接器的关键参数除接触电阻（优良的射频接头应在
毫欧量级）、绝缘电阻（通常要求大于
$1000\ \mathrm{M}\Omega$）与插拔次数（BNC 约 $500$ 次，
SMA 约 $500$ 次，板对板连接器往往仅数十次）之外，
还包括额定电流与电压、耐压强度、接触件镀层（镀金
接触电阻小且耐腐蚀，镀锡成本低但易氧化，两者不可
混配以免电偶腐蚀）、工作温度范围、防护等级与互调
指标。实际维修中，接头的氧化、进水与中心针退缩是
射频链路驻波异常与信号时断时续的常见原因，检查时
应先目视中心针位置与介质是否变形，再用万用表测
通断与绝缘，必要时用驻波表或网络分析仪定位。

\section{线材选型与应用}
\label{sec:wiring_selection}

选型时需考虑：

\begin{itemize}
    \item \textbf{电流容量}：根据载流量选择线径
    \item \textbf{频率特性}：高频应用需考虑趋肤效应
    \item \textbf{环境条件}：温度、湿度、化学腐蚀等
    \item \textbf{机械强度}：抗拉强度、柔韧性要求
\end{itemize}

高频应用中，趋肤效应使电流集中在导体表面，需使用多股绞合线或镀银导线。

线材选型是一个多约束的折中过程，其一般步骤如下。首先根据
最大工作电流并考虑敷设方式与环境温度的降额系数确定最小
截面；随后按允许压降复核该截面，若压降超限则加大线径；
再按工作频率判断是否需要按传输线处理，若需要则确定特性
阻抗、允许衰减与屏蔽要求；最后校核绝缘的耐压与耐温等级、
外护套的耐候与阻燃性能，以及弯曲半径与抗拉强度是否满足
机械安装条件。对于承受反复弯折的场合还应关注弯折寿命，
选用高股数的柔软电缆并预留适当的弯曲余量。

\textbf{趋肤效应}：交流电流在导体中并非均匀分布，频率越高
越集中于表面，这一现象称为趋肤效应。电流密度自表面向内
按指数衰减，衰减到表面值 $1/e$ 处的深度称为趋肤深度：
\begin{equation}
\delta = \sqrt{\frac{2\rho}{\omega\mu}} = \sqrt{\frac{\rho}{\pi f \mu}}
\end{equation}
对铜导体，该式可简化为工程估算式
$\delta \approx 66/\sqrt{f}$（$\delta$ 以 mm 计，$f$ 以 Hz 计）。
据此可算出：$1\ \mathrm{kHz}$ 时趋肤深度约
$2.1\ \mathrm{mm}$，$1\ \mathrm{MHz}$ 时约
$66\ \mu\mathrm{m}$，$10\ \mathrm{MHz}$ 时约
$21\ \mu\mathrm{m}$，$100\ \mathrm{MHz}$ 时仅约
$6.6\ \mu\mathrm{m}$。可见在射频频段电流几乎只在导体表层
极薄的一层中流动，导体的有效截面大为减小，交流电阻按
近似 $\sqrt{f}$ 的规律上升。由此可得三条实用结论：其一，
高频导体应增大表面积而非增大实心截面，故采用多股绞合线
（利兹线）或扁带；其二，导体表面的状态决定损耗，因而
镀银导线优于镀锡或裸铜，表面划伤与氧化都会增加损耗；
其三，管状导体与实心导体在高频下性能相当，故大功率
射频线圈与馈电杆多用铜管制成，既省材料又便于水冷。
此外，相邻导体间的邻近效应会进一步加剧电流分布不均，
故高频绕组应避免多层密绕。

\textbf{环境条件与机械强度}：温度直接决定绝缘寿命与载流
降额；湿度与凝露会降低绝缘电阻并加速接头腐蚀，户外
线缆须选用防水护套并在入户处设置滴水弯；紫外线会使
普通 PVC 护套粉化开裂，露天敷设应选用耐候等级的材料；
盐雾与化学腐蚀环境须选用镀层可靠的连接器并做密封处理；
存在鼠害或机械损伤风险的场合应采用铠装电缆或穿管保护。
机械方面须核算抗拉强度，垂直敷设的长馈线应设置承力
钢索，不得由连接器承受电缆自重；每种电缆都规定了最小
弯曲半径（一般为外径的 $5\sim10$ 倍，同轴电缆常取
$10$ 倍以上），过度弯折会使同轴结构的尺寸比 $D/d$
改变，导致局部阻抗失配与驻波恶化。

\textbf{施工与维护要点}：布线时应把强电线路与弱电信号
线路分开走线，必要时垂直交叉而非平行贴近，以减小容性
与感性耦合；射频馈线与音频线不可捆扎在一起。屏蔽层
接地要遵循前述原则，并保证接地路径短而粗。焊接接头
应做到饱满光亮，避免虚焊与冷焊，焊后清除助焊剂残留；
压接接头须使用配套压钳并做拉力抽检。所有接头与端子
应做标识，成束线扎不宜过紧以免压伤绝缘。日常维护中
应定期检查线缆有无破损与老化、接头有无发热变色与
氧化、驻波比与衰减是否较基线值劣化。把线材与连接器
视为与有源器件同等重要的组成部分，是保证通信设备
长期稳定运行的基本要求。
